\documentclass[11pt,a4paper]{article}
\usepackage[utf8]{inputenc}
\usepackage[T1]{fontenc}
\usepackage[french]{babel}
\usepackage{geometry}
\geometry{top=2.5cm, bottom=2.5cm, left=2.5cm, right=2.5cm}
\usepackage{amsmath,amssymb,amsfonts,mathtools}
\usepackage{xcolor}
\usepackage{fancyhdr}
\usepackage{booktabs}
\usepackage{caption}
\usepackage{titlesec}
\usepackage{graphicx}
\usepackage{hyperref}

% --- Palette de Couleurs Académique ---
\definecolor{mainblue}{RGB}{26, 54, 93}       % Bleu nuit pour les titres principaux
\definecolor{slate}{RGB}{74, 85, 104}         % Gris ardoise pour les sous-titres et détails
\definecolor{lightbg}{RGB}{247, 250, 252}     % Fond très clair pour les encadrés
\definecolor{bordergrey}{RGB}{226, 232, 240}  % Gris clair pour les lignes de séparation

% --- Configuration des Titres ---
\titleformat{\section}{\color{mainblue}\large\bfseries}{\thesection.}{0.5em}{}[{\titlerule[0.5pt]}]
\titleformat{\subsection}{\color{slate}\normalsize\bfseries}{\thesubsection.}{0.5em}{}

% --- Configuration des En-têtes et Pieds de page ---
\pagestyle{fancy}
\fancyhf{}
\renewcommand{\headrulewidth}{0.4pt}
\renewcommand{\footrulewidth}{0.4pt}
\fancyhead[L]{\footnotesize\color{slate}\textbf{Note Technique} | LAAS-CNRS}
\fancyhead[R]{\footnotesize\color{slate}Décomposition de Benders Robuste}
\fancyfoot[C]{\footnotesize\color{slate}Page \thepage}

% --- Configuration de Hyperref ---
\hypersetup{
    colorlinks=true,
    linkcolor=mainblue,
    urlcolor=mainblue,
    pdftitle={deroulement_algo_benders},
    pdfauthor={Mathis Francine-Habas}
}

\begin{document}

\begin{center}
    {\color{mainblue}\LARGE\bfseries Déroulement Manuel de l'Algorithme de Benders}\\[0.3cm]
    {\color{slate}\large Planification Robuste sous Incertitude}\\[0.5cm]
    \noindent\makebox[\linewidth]{\color{bordergrey}\rule{\linewidth}{1pt}}\\[0.2cm]
    \begin{minipage}{0.48\textwidth}
        \begin{flushleft} \small
            \textbf{Auteur :} Mathis Francine-Habas \\
            \textbf{Structure :} LAAS-CNRS
        \end{flushleft}
    \end{minipage}
    \hfill
    \begin{minipage}{0.48\textwidth}
        \begin{flushright} \small
            \textbf{Date :} \today \\
            \textbf{Version :} 1.0
        \end{flushright}
    \end{minipage}\\[0.1cm]
    \noindent\makebox[\linewidth]{\color{bordergrey}\rule{\linewidth}{1pt}}
\end{center}

\vspace{0.4cm}

% --- Section 1 : Cadre Mathématique ---
\section{Formulation et Configuration du Problème}
\label{sec:formulation}
\textit{Utilisez cette section pour poser brièvement le problème mathématique global (généralement un problème de type Min-Max ou à deux étapes sous incertitude) et sa partition en Problème Maître et Sous-Problème.}

\subsection{Problème Maître de Benders (PM)}
Le problème maître relaxé à l'itération $k$ s'exprime de la forme suivante :
\begin{equation}
    \begin{aligned}
        \min_{x, \theta} \quad & c^T x + \theta \\
        \text{s.t.} \quad & A x \ge b \\
        & \text{[Coupes de faisabilité générées]} \\
        & \text{[Coupes d'optimalité générées]} \\
        & x \in X, \, \theta \in \mathbb{R}
    \end{aligned}
\end{equation}

\subsection{Sous-Problème de Benders (SP)}
Pour une solution maître fixée $\hat{x}^{(k)}$ et un scénario d'incertitude donné (ou le pire scénario dans le cas robuste), le dual du sous-problème s'écrit :
\begin{equation}
    \begin{aligned}
        \max_{\lambda} \quad & (d - G\hat{x}^{(k)})^T \lambda \\
        \text{s.t.} \quad & W^T \lambda \le h \\
        & \lambda \ge 0
    \end{aligned}
\end{equation}

\vspace{0.3cm}

% --- Section 2 : Tableau Synthétique ---
\section{Tableau de Synthèse du Déroulement Étape par Étape}
\label{sec:tableau_synthese}
\textit{Ce tableau permet de suivre l'évolution globale des variables, des bornes et des coupes générées à chaque itération. C'est l'élément central pour valider la convergence de votre calcul à la main.}

\begin{table}[h!]
    \centering
    \caption{Évolution des indicateurs et solutions au fil des itérations}
    \label{tab:benders_iterations}
    \small
    \begin{tabular}{cccccccc}
        \toprule
        \textbf{Itér. ($k$)} & $\hat{x}^{(k)}$ & $\theta^{(k)}$ & \textbf{Statut SP} & $\lambda^{(k)}$ (ou rayon) & \textbf{Type de Coupe} & \textbf{Borne Inf ($LB$)} & \textbf{Borne Sup ($UB$)} \\
        \midrule
        1 & $x^{(1)} = \dots$ & $-\infty$ & Optimale / Infaisable & $\lambda^{(1)} = \dots$ & Optimalité / Faisabilité & $\dots$ & $\dots$ \\
        2 & $x^{(2)} = \dots$ & $\dots$ & $\dots$ & $\dots$ & $\dots$ & $\dots$ & $\dots$ \\
        3 & $\dots$ & $\dots$ & $\dots$ & $\dots$ & $\dots$ & $\dots$ & $\dots$ \\
        \bottomrule
    \end{tabular}
\end{table}

\pagebreak % Permet de passer proprement à l'analyse détaillée sur la page suivante

% --- Section 3 : Analyse Détaillée des Itérations ---
\section{Analyse Détaillée et Évolution des Éléments}
\label{sec:analyse_detaillee}
\textit{Pour chaque itération, décrivez précisément le changement des variables, la construction de la coupe et l'explication du choix de l'algorithme (pourquoi la borne bouge, pourquoi tel scénario est activé, etc.).}

\subsection{Itération 1}
\textbf{1. Résolution du Problème Maître (PM) :} \\
À l'itération 1, aucune coupe n'est encore présente. On obtient la solution initiale :
$\hat{x}^{(1)} = \dots$ et $\theta^{(1)} = -\infty$. Borne Inférieure $LB = \dots$

\noindent \textbf{2. Évaluation du Sous-Problème (SP) :} \\
En fixant $x = \hat{x}^{(1)}$, on résout le dual du sous-problème. 
\begin{itemize}
    \item \textbf{Résultat :} [Indiquer si le problème est optimal ou infaisable (génère un rayon extrême)].
    \item \textbf{Solution duale trouvée :} $\lambda^{(1)} = \dots$
\end{itemize}

\noindent \textbf{3. Génération de la Coupe et Mise à jour :} \\
La condition de convergence n'est pas vérifiée car $UB \neq LB$. On génère une \textbf{coupe d'optimalité/faisabilité} définie par :
\begin{equation*}
    \theta \ge (d - Gx)^T \lambda^{(1)} \implies \dots
\end{equation*}
\textit{Pourquoi cet élément change ?} La coupe vient restreindre l'espace de recherche du Problème Maître pour forcer $\theta$ (ou la faisabilité) à s'ajuster lors de l'étape suivante.

\subsection{Itération 2}
\textbf{1. Résolution du Problème Maître (PM) :} \\
On réintègre la coupe précédente. La nouvelle solution du PM donne $\hat{x}^{(2)} = \dots$ et $\theta^{(2)} = \dots$. La borne inférieure remonte à $LB = \dots$.

\noindent \textbf{2. Évaluation du Sous-Problème (SP) :} \\
Avec $\hat{x}^{(2)}$, le sous-problème donne \dots

\vspace{0.6cm}

% --- Section 4 : Visualisation / Figures ---
\section{Visualisation Graphique de la Convergence}
\label{sec:visualisation}
\textit{Insérez ici vos graphiques montrant le resserrement des bornes (Convergence Plot) ou l'arborescence des coupes dans l'espace des solutions.}

\begin{figure}[h!]
    \centering
    % \includegraphics[width=0.75\linewidth]{votre_graphique_convergence.png} % Décommenter une fois l'image prête
    \begin{minipage}{0.75\linewidth}
        \centering
        \color{slate}
        \framebox[\linewidth]{\raisebox{0pt}[3cm][3cm]{\dots Encadré de placement pour votre Figure (Bornes Sup/Inf en fonction des Itérations) \dots}}
    \end{minipage}
    \caption{Graphique de convergence : Évolution de la borne inférieure ($LB$) et de la borne supérieure ($UB$)}
    \label{fig:convergence}
\end{figure}

\end{document}